% CV MBA Analytique affaires - Stephane KPOVIESSI
% Candidature au MBA analytique affaires - Universite Laval - automne 2027
% Base : CV Finance (Universite de Sherbrooke), repositionne pour une candidature universitaire en strategie.
% Compile target: pdfLaTeX

\documentclass[10pt,a4paper]{article}

\usepackage[a4paper,left=0.95cm,right=0.95cm,top=0.55cm,bottom=0.55cm]{geometry}
\usepackage[utf8]{inputenc}
\usepackage[T1]{fontenc}
\usepackage{lmodern}
\usepackage{xcolor}
\usepackage{enumitem}
\usepackage[hidelinks]{hyperref}
\usepackage{tabularx}
\usepackage{titlesec}
\usepackage{fontawesome5}

\definecolor{accent}{HTML}{1F4E79}
\definecolor{accentdark}{HTML}{163A5A}
\definecolor{text}{HTML}{2F3437}
\definecolor{muted}{HTML}{666D70}
\definecolor{rule}{HTML}{D5E2EC}
\definecolor{sidebg}{HTML}{F2F7FB}

\pagestyle{empty}
\setlength{\parindent}{0pt}
\setlength{\parskip}{3.2pt}
\setlist[itemize]{leftmargin=*, topsep=2pt, itemsep=2pt, parsep=0pt, partopsep=0pt}

\titleformat{\section}
  {\fontsize{11.2pt}{13pt}\selectfont\bfseries\color{accent}}
  {}{0pt}{}
  [\vspace{2pt}{\color{rule}\titlerule}\vspace{3pt}]
\titlespacing*{\section}{0pt}{8pt}{4pt}

\newcommand{\smallsection}[2]{%
  \vspace{7pt}
  {\fontsize{8.9pt}{10.45pt}\selectfont\bfseries\color{accentdark}%
  {\scriptsize\color{accent}#1}\hspace{3.5pt}#2}\par
  \vspace{3pt}
}

\newcommand{\contactline}[2]{%
  \noindent\makebox[0.34cm][l]{{\color{accent}\scriptsize#1}}%
  \begin{minipage}[t]{\dimexpr\linewidth-0.36cm\relax}%
  \fontsize{8.15pt}{9.5pt}\selectfont #2%
  \end{minipage}\par
  \vspace{1pt}
}

\newcommand{\sideitem}[2]{%
  \textbf{#1}\par
  #2\par
  \vspace{5pt}
}

\newcommand{\certitem}[2]{%
  \textbf{#1}\par
  {\fontsize{7.15pt}{8.1pt}\selectfont\color{muted}#2}\par
  \vspace{4pt}
}

\newcommand{\role}[3]{%
  \vspace{2pt}
  \begin{tabularx}{\linewidth}{@{}Xr@{}}
  \textbf{#1} & {\color{muted}#2}
  \end{tabularx}\vspace{0.8pt}\\
  {\color{muted}#3}\par
  \vspace{2pt}
}

\newsavebox{\maincolumnbox}
\newlength{\maincolumnheight}
\newlength{\sideboxheight}

\begin{document}
\fontsize{10pt}{12.3pt}\selectfont
\sloppy

\begin{minipage}[t]{0.52\textwidth}
{\LARGE\bfseries\color{accentdark} Stéphane KPOVIESSI}\par
\vspace{1pt}
{\large\color{accent}Ingénieur diplômé en Big Data \& Data Science}\par
\vspace{1pt}
{\color{muted}Candidature au MBA -- Analytique d\textquotesingle affaires\\ Université Laval \textbullet{} Automne 2027}
\end{minipage}
\hfill
\begin{minipage}[t]{0.42\textwidth}
\raggedleft
\vspace{2pt}
{\normalsize\bfseries\color{accent}Analyser les données \textbullet{} Éclairer les décisions}\par
\vspace{3pt}
{\small\color{muted}
Relier l'analyse des données aux enjeux de gestion, de finance et de performance.
}
\end{minipage}

\vspace{5pt}
{\color{rule}\hrule height 1pt}
\vspace{6pt}

\begin{lrbox}{\maincolumnbox}
\begin{minipage}[t]{0.695\textwidth}
\vspace{0pt}
\raggedright

\section{Profil}
Ingénieur diplômé en Big Data et Data Science, avec des expériences en fiabilisation d'indicateurs chez Allianz France, en analyse d'activité chez Lyne \& Frères et en R\&D chez Sylvagreg. Je souhaite compléter mon expertise technique par une formation en gestion, en stratégie et en finance pour conduire des projets d'analytique d'affaires et accompagner les décisions des organisations.

\section{Expériences}
\role{Seshat - Fondateur \& Ingénieur IA}{Sept. 2026 - Aujourd'hui}{Projet entrepreneurial en construction - IA, automatisation et logiciels métiers}
\begin{itemize}
  \item Définition de l'offre visée (solutions d'IA, d'automatisation et de logiciels métiers pour les entreprises) et conduite du développement produit jusqu'à la mise en ligne.
  \item Conception et développement de la plateforme (Python/Go, application desktop Electron/React) : ingestion documentaire, RAG hybride et orchestration d'agents.
  \item Analyse comparative des fournisseurs et des briques techniques (modèles de langage, recherche vectorielle, hébergement) pour arbitrer les dépendances du produit.
\end{itemize}
{\footnotesize\color{accent}\faIcon{globe}\ \href{https://seshat-ai.com}{seshat-ai.com}\enspace\textbullet\enspace\faIcon{github}\ \href{https://github.com/KPO-Tech/seshat}{github.com/KPO-Tech/seshat}}\par

\vspace{4pt}
\role{Allianz France - Ingénieur Data / Data Scientist}{Nov. 2025 - Avr. 2026}{Direction P4D, Advanced Data \& Climate - Paris}
\begin{itemize}
  \item Migration du datamart Assurance de SAS vers Python/PySpark : cadrage des traitements avec les équipes métier, validation de la parité fonctionnelle sur plus de 15 000 lignes de code et sécurisation de la reprise des indicateurs de pilotage.
  \item Architecture Médaillon sur Azure Databricks intégrant plus de 150 règles de gestion métier de l'assurance : des indicateurs reproductibles et auditables.
  \item Contrôle qualité des données et tests automatisés : détection des écarts avant restitution et fiabilisation des indicateurs métier.
\end{itemize}

\vspace{4pt}
\role{LYNE ET FRERES SARL - Analyste de données}{2023 - 2025}{Appui à la décision - Bénin / À distance}
\begin{itemize}
  \item Analyses commerciales, opérationnelles et agricoles (Python/Pandas) : suivi de l'activité, analyse des écarts et mise en évidence des points de tension pour la direction.
  \item Rapports et tableaux de bord (Streamlit / web) restituant les analyses et les indicateurs auprès des décideurs.
  \item Automatisation de traitements et outils internes pour homogénéiser les chiffres utilisés par les équipes métier.
\end{itemize}

\vspace{4pt}
\role{Groupe Sylvagreg - Machine Learning Engineer}{Mai 2024 - Août 2024}{Lille, France}
R\&D en vision par ordinateur : reconstruction 3D à partir d'images 2D (PyTorch), cadrage du besoin avec les équipes, préparation des données, expérimentation et déploiement (Docker, AWS).\par

\section{Engagements \& gouvernance}
\textbf{JUNIA ALUMNI} - association des diplômés de JUNIA\\
\begin{itemize}
  \item \textbf{Vice-président (2024 - 2025)} : coordination du bureau et des bénévoles, préparation des décisions associatives, relations avec l'école et les partenaires diplômés.
  \item \textbf{Responsable communication (2023 - 2024)} : animation de la communauté, supports et canaux de diffusion, organisation d'événements alumni.
\end{itemize}

\end{minipage}
\end{lrbox}
% Le bandeau est dimensionne sur la colonne de droite, plus courte : on ajoute une marge basse
% pour que la derniere ligne du bandeau reste dans le fond bleu.
\setlength{\maincolumnheight}{\dimexpr\ht\maincolumnbox+\dp\maincolumnbox\relax}
\setlength{\sideboxheight}{\dimexpr\maincolumnheight+14pt\relax}

\begin{minipage}[t]{0.285\textwidth}
\vspace{0pt}

\hspace*{-0.18cm}%
\colorbox{sidebg}{%
  \hspace*{0.18cm}%
  \begin{minipage}[t]{0.91\linewidth}
  \vspace{8pt}
  \fontsize{8.45pt}{10.35pt}\selectfont
  \raggedright

  \smallsection{\faIcon{address-card}}{Contact}
  \contactline{\faIcon{map-marker-alt}}{Cergy / Île-de-France}
  \contactline{\faIcon{phone}}{+33 7 48 59 47 13}
  \contactline{\faIcon{envelope}}{\href{mailto:oastephanekpoviessi@gmail.com}{oastephanekpoviessi@gmail.com}}
  \contactline{\faIcon{link}}{\href{https://linkedin.com/in/stephanekpoviessi}{linkedin.com/in/stephanekpoviessi}}
  \contactline{\faIcon{github}}{\href{https://github.com/KPO-Tech}{github.com/KPO-Tech}}
  \contactline{\faIcon{globe}}{\href{https://seshat-ai.com}{seshat-ai.com}}

  \smallsection{\faIcon{graduation-cap}}{Formation}
  \sideitem{Un seul cursus JUNIA}{Cinq années d'études, 2020--2025. Diplôme d'ingénieur ISEN ; jury du 18/09/2025, attestation du 20/04/2026.}
  \sideitem{Lille, France (2023--2025)}{4\textsuperscript{e} et 5\textsuperscript{e} années -- Big Data \& Data Science. Visualisation : 17,50/20 ; Hadoop \& Spark : 17,50/20.}
  \sideitem{Rabat, Maroc (2022--2023)}{3\textsuperscript{e} année -- Électronique et numérique. Aide à la décision : 18,50/20.}
  \sideitem{Cotonou, Bénin (2020--2022)}{1\textsuperscript{re} et 2\textsuperscript{e} années -- Prépa Ste Marie-Stella. Moyennes : 16,00 ; 16,42 ; 16,00 ; 16,94/20. En 2\textsuperscript{e} année S2 : maths 19,70/20 et informatique 19,45/20.}

  \smallsection{\faIcon{layer-group}}{Compétences}
  \sideitem{Diagnostic \& analyse}{Formuler une problématique, analyser des données et des processus, mesurer un écart.}
  \sideitem{Appui à la décision}{Indicateurs de pilotage, restitution aux décideurs, qualité et fiabilité des chiffres.}
  \sideitem{Coordination}{Cadrage avec les équipes métier, travail en transverse, accompagnement à la prise en main des outils.}
  \sideitem{Données \& outils}{Python, SQL, Pandas, PySpark, PostgreSQL, Azure Databricks, Docker, Git.}

  \smallsection{\faIcon{language}}{Langues}
  \textbf{Français} - langue maternelle\par
  \textbf{Anglais} - B2, Linguaskill Business\par

  \smallsection{\faIcon{certificate}}{Certifications}
  \certitem{\href{https://www.coursera.org/account/accomplishments/verify/MWJS4Q3S2C3L}{Supervised ML: Regression \& Classification}}{DeepLearning.AI / Stanford Online \textbullet{} 2023}
  \certitem{\href{https://www.freecodecamp.org/certification/amiche/data-analysis-with-python-v7}{Data Analysis with Python}}{freeCodeCamp \textbullet{} 2023}


  \smallsection{\faIcon{users}}{Champs d'intérêt}
  {\color{muted}Stratégie \textbullet{} finance d\textquotesingle entreprise \textbullet{} pilotage de la performance \textbullet{} entrepreneuriat.}\par
  \vspace{3pt}
  \end{minipage}
}

\end{minipage}
\hfill
\usebox{\maincolumnbox}

\end{document}
